% Part: computability
% Chapter: computability-theory
% Section: introduction

\documentclass[../../../include/open-logic-section]{subfiles}

\begin{document}

\olfileid{cmp}{thy}{int}
\olsection{પરિચય}

તર્કશાસ્ત્રની \emph{સંગણનીયતા સિદ્ધાંત} નામની શાખા વિધેયો અને ગણોની
સંગણનીયતા અથવા સાપેક્ષ સંગણનીયતા અંગેના પ્રશ્નોનો અભ્યાસ કરે છે.
આ વિષય પર ક્લીનીનો પ્રભાવ એવો રહ્યો છે કે પહેલાં તેને \emph{પુનરાવર્તન
  સિદ્ધાંત} કહેવામાં આવતો હતો; આજે બંને નામ સામાન્ય રીતે વપરાય છે.

જો સંગણનાના કોઈ નિદર્શમાં $f\colon \Nat \pto \Nat$ પ્રકારનું વિધેય
ગણી શકાય, તો તેને \emph{આંશિક સંગણનીય} કહીએ. $f$ સંપૂર્ણ હોય તો
$f$ને માત્ર \emph{સંગણનીય} કહીએ. જે સંબંધ~$R$નું લક્ષણ વિધેય
$\Char{R}$ સંગણનીય હોય તેને પણ સંગણનીય કહીએ. $f$ અને~$g$ આંશિક
વિધેયો હોય, ત્યારે $f(x) \fdefined$નો અર્થ $x$ પર $f$ વ્યાખ્યાયિત
છે, એટલે $x$ $f$ના પ્રદેશમાં છે, એવો થશે; અને $f(x) \fundefined$નો
અર્થ $f$ $x$ પર વ્યાખ્યાયિત નથી એવો થશે. $f(x) \simeq g(x)$થી એવો અર્થ લઈશું કે
$f(x)$ અને $g(x)$ બંને અનિર્ધારિત છે, અથવા બંને વ્યાખ્યાયિત અને
સમાન છે.

સંગણનાના કોઈ એક ચોક્કસ નિદર્શનો ઉલ્લેખ કર્યા વિના પણ સંગણનીયતાનો
સિદ્ધાંત તપાસી શકાય છે. તે માટે સર્વવ્યાપી આંશિક સંગણનીય વિધેય
$\fn{Un}(k,x)$ છે એવું બતાવાય છે. તેના વડે આંશિક સંગણનીય વિધેયોની
પરિગણના કરી શકાય છે. $k$મા એક-આર્ગ્યુમેન્ટવાળા આંશિક સંગણનીય
વિધેય માટે $\cfind{k}$ સંકેત વાપરીશું, જ્યાં
$\cfind{k}(x) \simeq \fn{Un}(k,x)$ છે. (ક્લીનીએ આ માટે
$\{ k \}$ વાપર્યું હતું, પરંતુ તાજેતરમાં આ સંકેત ઓછો વપરાય છે.)
થોડું વધુ સામાન્ય રીતે, કોઈ પણ આર્ગ્યુમેન્ટ-સંખ્યાવાળાં આંશિક
સંગણનીય વિધેયોની એકસરખી રીતથી પરિગણના કરી શકાય છે. $k$મા
$n$-આર્ગ્યુમેન્ટવાળા આંશિક પુનરાવર્તી વિધેય માટે
$\cfind{k}[n]$ લખીશું.

$f(\vec x,y)$ સંપૂર્ણ અથવા આંશિક વિધેય હોય, તો
$\umin{y}{f(\vec x,y)}$ એ $\vec x$નું એવું વિધેય છે જે સૌથી નાનું
$y$ આપે છે જેના માટે $f(\vec x,y)=0$ છે, શરતે કે
$f(\vec x,0)$, \dots, $f(\vec x,y-1)$ બધાં વ્યાખ્યાયિત હોય.
આવું $y$ ન હોય તો $\umin{y}{f(\vec x,y)}$ અનિર્ધારિત છે.
$R(\vec x,y)$ સંબંધ હોય, તો $\umin{y}{R(\vec x,y)}$ એ
$R(\vec x,y)$ સાચું હોય એવું સૌથી નાનું~$y$ છે; બીજા શબ્દોમાં,
$1 \tsub \Char{R}(\vec x,y)=0$ હોય એવું સૌથી નાનું~$y$.

વિધેય સંગણનીય છે એવું બતાવવા બે રીતે આગળ વધી શકાય:
\begin{enumerate}
\item ચુસ્ત રીતે: ટ્યુરિંગ મશીન અથવા આંશિક પુનરાવર્તી વિધેય
  સ્પષ્ટ રીતે વર્ણવીને તે ઇચ્છિત વિધેય ગણે છે એવું બતાવવું;
\item અનૌપચારિક રીતે: વિધેય ગણતી કલનવિધિ વર્ણવવી અને ચર્ચની
  માન્યતાનો આધાર લેવો.
\end{enumerate}
આ બે રીતો વચ્ચે તીક્ષ્ણ સીમા નથી. કલનવિધિનું વિગતવાર વર્ણન એવું
માહિતી આપે કે સિદ્ધાંતમાં તેના માટે યોગ્ય ટ્યુરિંગ મશીન કે આંશિક
પુનરાવર્તી વ્યાખ્યાઓની શ્રેણી કેવી રીતે ઘડી શકાય તે પ્રમાણમાં સ્પષ્ટ
થવું જોઈએ. સંપૂર્ણ ઔપચારિક વ્યાખ્યાઓ ઘણી વાર સમજમાં ખાસ વધારો
કરતી નથી. તેથી બંને રીતો વચ્ચે યોગ્ય સમતુલા રાખીશું: ચોક્કસ
વ્યાખ્યાઓ મેળવી શકાય એવી ખાતરી આપતી, છતાં સહજ સમજમાં ઊતરે એવી
ચુસ્ત સાબિતીઓ શોધીશું.

\end{document}
